\section{Application Architecture}\label{sec:studio-architecture}

Motivo Studio has three runtime services in the Docker Compose stack.
The client is a Next.js application with React components for the editor, logs, piano-roll visualization, playback, and
theme controls~\cite{NextJS_2026,React_2026}.
The server is an Express API that validates compile requests with Zod and invokes the native \texttt{motivoc}
binary~\cite{Express_2026,Zod_2026}.
The gateway is an nginx container that exposes the composed application through HTTPS at
\texttt{motivo-studio.local}.
This separation keeps the C++ compiler outside the browser while still giving the frontend an interactive compile loop.

\begin{figure}[htbp]
  \centering
  \begin{tikzpicture}[
      service/.style={rounded corners=5pt, draw=motivoInk!65, very thick, minimum width=4.0cm, minimum height=1.25cm,
      align=center, font=\small\sffamily},
      file/.style={rounded corners=3pt, draw=motivoInk!45, thick, minimum width=2.9cm, minimum height=0.75cm,
      align=center, font=\scriptsize\sffamily},
      arrow/.style={-Stealth, very thick, draw=motivoInk!75}
    ]
    \node[service, fill=motivoBlue!14] (client) {Next.js client\\Monaco editor and playback};
    \node[service, fill=motivoCyan!18, right=1.1cm of client] (server) {Express server\\compile API};
    \node[service, fill=motivoGreen!18, right=1.1cm of server] (compiler) {Native compiler\\\texttt{motivoc}};

    \node[file, fill=motivoAmber!18, below=0.65cm of client] (midi) {MIDI bytes};
    \node[file, fill=motivoRose!12, below=0.65cm of server] (diag) {diagnostics};
    \node[file, fill=motivoSoft, below=0.65cm of compiler] (midfile) {\texttt{.mid} file};

    \draw[arrow] (client) -- node[above, font=\scriptsize\sffamily] {\texttt{POST /compile}} (server);
    \draw[arrow] (server) -- node[above, font=\scriptsize\sffamily] {temporary source} (compiler);
    \draw[arrow] (compiler) -- (midfile);
    \draw[arrow] (midfile) -- (server);
    \draw[arrow] (server) -- node[below, font=\scriptsize\sffamily] {binary or errors} (client);
    \draw[arrow] (server) -- (diag);
    \draw[arrow] (client) -- (midi);
  \end{tikzpicture}
  \caption{Motivo Studio runtime architecture.
    The browser does not compile C++ code directly; it delegates compilation to the
  server.}\label{fig:studio-architecture}
\end{figure}

The frontend is organized by feature directories.
The editor feature wraps Monaco, registers Motivo language tokens, defines light and dark themes, and maps compiler
diagnostics into Monaco markers.
The compile feature owns the HTTP client and the response types.
The MIDI feature parses generated binary data into browser-side structures.
The piano-roll feature renders note rectangles on a canvas.
The playback feature schedules browser audio and exposes controls for play, stop, progress, and MIDI download.
This feature-oriented layout keeps UI responsibilities closer to their domain than a flat component directory would.

While the core contribution of this thesis is the compiler, a web-based \textit{Playground} was developed to demonstrate
the language's capabilities and provide an intuitive environment for composition.
The playground integrates the C++ compiler into a modern full-stack web application.

\subsection{System Architecture}\label{detail-subsec:system-architecture}

The playground follows a distributed architecture:

\noindent \textbf{Frontend (Next.js).} A React-based interface that provides the editor and visualization tools.

\noindent \textbf{Backend (Express).} A Node.js API that acts as a bridge.
It receives Motivo source code from the client, invokes the C++ \texttt{motivoc} binary via a sub-process, and streams
the resulting MIDI file back to the browser.

\noindent \textbf{Integration.} The system is containerized using Docker, ensuring that the C++ environment (including
Bison/Flex runtimes) is identical across development and production.

The backend is smaller but equally important.
The Express app exposes \texttt{GET /health} and \texttt{POST /compile}.
The compile route validates the request body with Zod, writes the source code to a temporary \texttt{.motivo} file,
spawns \texttt{motivoc}, and returns either \texttt{audio/midi} bytes or structured diagnostics.
The diagnostics parser normalizes compiler output into severity, stage, message, and optional source location fields so
the browser can highlight the relevant source range.
